\chapter{\texorpdfstring{$W_2$}{W2} と Gaussian 測度}
\label{ch:gaussian}

\section{設定と方針}
\label{sec:gauss-setting}

本章では $X=\R^n$（Euclid 距離 $d(x,y)=\norm{x-y}$），$p=2$ とし，
本資料の最終目標である閉形式
\[
  \Wass_2\bigl(\Normal(m_1,\Sigma_1),\,\Normal(m_2,\Sigma_2)\bigr)^2
  =
  \norm{m_1-m_2}^2
  +\tr\Bigl(\Sigma_1+\Sigma_2
    -2\bigl(\Sigma_1^{1/2}\Sigma_2\Sigma_1^{1/2}\bigr)^{1/2}\Bigr)
\]
（Givens--Shortt (1984) Proposition 7）を証明する．
証明は次の4段階からなる．
\begin{enumerate}
  \item \textbf{平均の分離}：$\Wass_2^2$ は
    「平均のずれの2乗」と「中心化した測度間の $\Wass_2^2$」に分かれる．
  \item \textbf{コストの行列表示}：中心化後，coupling $\pi$ のコストは
    相互共分散行列 $K_\pi$ を使って
    $\tr\Sigma_1+\tr\Sigma_2-2\tr K_\pi$ と書け，
    ブロック行列 $\bigl(\begin{smallmatrix}\Sigma_1&K_\pi\\K_\pi^\top&\Sigma_2\end{smallmatrix}\bigr)$
    は半正定値になる．
  \item \textbf{トレース最大化}：この制約下での $\tr K$ の最大値は
    $\tr\bigl(\Sigma_1^{1/2}\Sigma_2\Sigma_1^{1/2}\bigr)^{1/2}$
    である（純粋な線形代数）．
    これで下界が出る．しかもこの下界は Gaussian に限らず，
    同じ平均・共分散をもつ任意の測度で成り立つ．
  \item \textbf{達成}：最大化する $K^*$ を共分散にもつ
    結合 Gaussian 測度が，下界を達成する coupling になる．
\end{enumerate}
原論文は共分散行列が正則な場合のみを扱い
（特異な場合は「同様だが煩雑な場合分けを要する」と注記のみ：p.~239），
証明には最適 coupling の存在（弱コンパクト性）と
Lagrange 乗数法・Courant--Fischer の定理を用いる．
本章では，逆行列を使わないブロック半正定値行列の特徴づけ
（補題~\ref{lem:la-block-psd}）と結合 Gaussian の明示的構成により，
\textbf{特異な共分散も込めた一般の場合}を，
弱収束の理論を使わずに証明する．

\begin{remark}{$\R^n$ は Polish 空間}{gauss-rn-polish}
  $(\R^n,\norm{\cdot})$ は完備かつ可分（定義~\ref{def:meas-polish}）である．
  実際，$\R^n$ の Cauchy 列は各座標が $\R$ の Cauchy 列であり，
  実数の完備性（既知とする）から各座標が収束するので全体も収束する．
  また $\mathbb{Q}^n$ は可算で，各座標を有理数で近似すれば稠密である．
  したがって第\ref{ch:wasserstein-metrics}章の理論がそのまま使える．
  以下，基点を $x_0=0$ に取り（補題~\ref{lem:wp-p-indep}），
  \[
    \Pp{2}(\R^n)
    =\Bigl\{\mu\in\Prob(\R^n)\;\Bigm|\;\int_{\R^n}\norm{x}^2\,\d\mu(x)<\infty\Bigr\}
  \]
  とする．
\end{remark}

\section{平均と共分散}
\label{sec:gauss-mean-cov}

\begin{definition}{平均ベクトルと共分散行列}{gauss-mean-cov}
  $\mu\in\Pp{2}(\R^n)$ に対し，\textbf{平均ベクトル}と\textbf{共分散行列}を
  \[
    m_\mu:=\int_{\R^n}x\,\d\mu(x)\in\R^n,
    \qquad
    \Sigma_\mu:=\int_{\R^n}(x-m_\mu)(x-m_\mu)^\top\,\d\mu(x)\in\R^{n\times n}
  \]
  と成分ごとの積分で定める．
  各成分は実際に可積分である：
  $\abs{x_i}\leq\frac{1}{2}(1+\norm{x}^2)$ より $m_\mu$ の各成分は有限，
  $\abs{(x-m_\mu)_i(x-m_\mu)_j}\leq\norm{x-m_\mu}^2
  \leq2\norm{x}^2+2\norm{m_\mu}^2$ より $\Sigma_\mu$ の各成分も有限である．
\end{definition}

\begin{claim}{平均・共分散の基本性質}{gauss-cov-basic}
  $\mu\in\Pp{2}(\R^n)$ とする．
  \begin{enumerate}
    \item $\Sigma_\mu$ は対称かつ半正定値で，任意の $u\in\R^n$ に対し
      $u^\top\Sigma_\mu u=\int\inner{u}{x-m_\mu}^2\,\d\mu$ である．
    \item $\displaystyle\int_{\R^n}\norm{x}^2\,\d\mu
      =\norm{m_\mu}^2+\tr\Sigma_\mu$．
      より一般に $L\in\R^{k\times n}$ に対し
      $\displaystyle\int_{\R^n}\norm{L(x-m_\mu)}^2\,\d\mu=\tr(L\Sigma_\mu L^\top)$．
    \item （アフィン像の平均・共分散）$L\in\R^{k\times n}$，$b\in\R^k$，
      $\ell(x):=Lx+b$ とすると，$\nu:=\ell\pushforward\mu\in\Pp{2}(\R^k)$ で
      \[
        m_\nu=Lm_\mu+b,
        \qquad
        \Sigma_\nu=L\Sigma_\mu L^\top.
      \]
  \end{enumerate}
\end{claim}

\begin{proof}
  (i)：対称性は定義から明らか．
  積分の線形性（成分ごと）より
  \[
    u^\top\Sigma_\mu u
    =\int u^\top(x-m_\mu)(x-m_\mu)^\top u\,\d\mu
    =\int\inner{u}{x-m_\mu}^2\,\d\mu\geq0.
  \]

  (ii)：後半から示す．$\norm{L(x-m_\mu)}^2=\tr\bigl(L(x-m_\mu)(x-m_\mu)^\top L^\top\bigr)$
  （主張~\ref{clm:la-trace-basic}(iii) を $u=v=L(x-m_\mu)$ に適用）であり，
  トレースと積分はともに成分の線形結合だから交換できて
  $\int\norm{L(x-m_\mu)}^2\,\d\mu=\tr(L\Sigma_\mu L^\top)$ である．
  前半は
  $\norm{x}^2=\norm{x-m_\mu}^2+2\inner{x-m_\mu}{m_\mu}+\norm{m_\mu}^2$
  を積分すると，中央の項は $\int(x-m_\mu)\,\d\mu=0$ より消え，
  $\int\norm{x-m_\mu}^2\,\d\mu=\tr\Sigma_\mu$（$L=I$ の場合）から従う．

  (iii)：まず
  $\norm{Lx+b}^2\leq2\norm{L}_{\mathrm{op}}^2\norm{x}^2+2\norm{b}^2$
  と変数変換公式（補題~\ref{lem:meas-change-of-variables}）から
  $\int\norm{y}^2\,\d\nu=\int\norm{Lx+b}^2\,\d\mu<\infty$，
  すなわち $\nu\in\Pp{2}(\R^k)$ である．
  平均は，変数変換公式（成分ごと）と積分の線形性より
  $m_\nu=\int(Lx+b)\,\d\mu=Lm_\mu+b$．
  共分散は $y-m_\nu=L(x-m_\mu)$（$y=Lx+b$）だから
  \[
    \Sigma_\nu
    =\int L(x-m_\mu)(x-m_\mu)^\top L^\top\,\d\mu
    =L\Sigma_\mu L^\top.  \qedhere
  \]
\end{proof}

\section{Gaussian 測度}
\label{sec:gauss-measure}

\begin{lemma}{Gauss 積分}{gauss-integral}
  $\varphi(t):=(2\pi)^{-1/2}e^{-t^2/2}$ とおくと，
  \[
    \int_\R\varphi\,\d\lambda_1=1,
    \qquad
    \int_\R t\,\varphi(t)\,\d\lambda_1(t)=0,
    \qquad
    \int_\R t^2\varphi(t)\,\d\lambda_1(t)=1.
  \]
\end{lemma}

\begin{proof}
  $I:=\int_\R e^{-t^2/2}\,\d\lambda_1$ とおく．
  Fubini--Tonelli の定理（定理~\ref{thm:meas-product}）と
  極座標変換（定理~\ref{thm:meas-polar}）から
  \[
    I^2
    =\int_{\R^2}e^{-(s^2+t^2)/2}\,\d\lambda_2
    =\int_0^\infty\!\!\int_0^{2\pi}e^{-r^2/2}\,r\,\d\theta\,\d r
    =2\pi\int_0^\infty re^{-r^2/2}\,\d r
    =2\pi\Bigl[-e^{-r^2/2}\Bigr]_0^\infty
    =2\pi
  \]
  であり，$I>0$ だから $I=\sqrt{2\pi}$，すなわち $\int\varphi\,\d\lambda_1=1$．

  次に $\int_0^\infty t\varphi(t)\,\d t=\varphi(0)<\infty$ だから
  $t\varphi(t)$ は可積分である．
  鏡映 $r(t):=-t$ は $\abs{\det(-1)}=1$ のアフィン写像なので
  $r\pushforward\lambda_1=\lambda_1$（定理~\ref{thm:meas-affine-cov}），
  変数変換公式（補題~\ref{lem:meas-change-of-variables}）と
  $\varphi(-t)=\varphi(t)$ から
  \[
    \int t\varphi(t)\,\d\lambda_1
    =\int(-t)\varphi(-t)\,\d\lambda_1
    =-\int t\varphi(t)\,\d\lambda_1,
  \]
  よって $\int t\varphi\,\d\lambda_1=0$ である．

  最後に，部分積分 $\bigl(-e^{-t^2/2}\bigr)'=te^{-t^2/2}$ より
  \[
    \int_{-R}^Rt^2e^{-t^2/2}\,\d t
    =\Bigl[-te^{-t^2/2}\Bigr]_{-R}^R+\int_{-R}^Re^{-t^2/2}\,\d t
    \;\xrightarrow{R\to\infty}\;0+\sqrt{2\pi}
  \]
  であり，左辺は非負関数の増大列の積分だから
  単調収束定理（定理~\ref{thm:meas-mct}）で
  $\int_\R t^2e^{-t^2/2}\,\d\lambda_1=\sqrt{2\pi}$，
  すなわち $\int t^2\varphi\,\d\lambda_1=1$ である．
\end{proof}

\begin{definition}{標準 Gaussian 測度}{gauss-standard}
  補題~\ref{lem:gauss-integral}より
  $\varphi_n(x):=(2\pi)^{-n/2}e^{-\norm{x}^2/2}$ は
  $\int_{\R^n}\varphi_n\,\d\lambda_n=1$ をみたす
  （Fubini--Tonelli の定理~\ref{thm:meas-product}で座標ごとに分解）．
  そこで密度をもつ測度（定義~\ref{def:meas-density}）として
  \[
    \gamma_n:=\varphi_n\,\d\lambda_n\in\Prob(\R^n)
  \]
  と定め，\textbf{標準 Gaussian 測度}とよぶ．
\end{definition}

\begin{claim}{$\gamma_n$ の直積構造とモーメント}{gauss-standard-moments}
  \begin{enumerate}
    \item $\gamma_n=\gamma_1\otimes\cdots\otimes\gamma_1$（$n$ 個の直積）であり，
      $1\leq r<n$ のとき，最初の $r$ 座標への射影
      $\mathrm{pr}_r(x_1,\dots,x_n):=(x_1,\dots,x_r)$ について
      ${\mathrm{pr}_r}\pushforward\gamma_n=\gamma_r$ である．
    \item $\gamma_n\in\Pp{2}(\R^n)$ で，平均 $0$，共分散 $I_n$，
      $\int\norm{x}^2\,\d\gamma_n=n$ である．
  \end{enumerate}
\end{claim}

\begin{proof}
  (i)：矩形集合 $A_1\times\cdots\times A_n$ 上で，
  密度の定義と Fubini--Tonelli の定理（定理~\ref{thm:meas-product}）から
  \[
    \gamma_n(A_1\times\cdots\times A_n)
    =\int\mathbf{1}_{A_1}(x_1)\varphi(x_1)\cdots\mathbf{1}_{A_n}(x_n)\varphi(x_n)
      \,\d\lambda_n
    =\prod_{i=1}^n\gamma_1(A_i)
  \]
  であり，これは $\gamma_1^{\otimes n}$ と一致する．
  矩形集合の全体は $\pi$-系で $\Borel(\R^n)$ を生成する
  （注意~\ref{rem:meas-borel-product}を繰り返し使う）から，
  一致定理（定理~\ref{thm:meas-uniqueness}）より
  $\gamma_n=\gamma_1^{\otimes n}$ である．
  同様に $\gamma_n=\gamma_r\otimes\gamma_{n-r}$ だから，
  直積測度の周辺分布（主張~\ref{clm:meas-product-marginals}）より
  \[
    ({\mathrm{pr}_r}\pushforward\gamma_n)(A)
    =\gamma_n(A\times\R^{n-r})
    =\gamma_r(A)
    \qquad(A\in\Borel(\R^r)).
  \]

  (ii)：非負関数への Fubini--Tonelli の定理と
  補題~\ref{lem:gauss-integral}から
  $\int\norm{x}^2\,\d\gamma_n=\sum_{i=1}^n\int x_i^2\,\d\gamma_n=n<\infty$，
  特に $\gamma_n\in\Pp{2}(\R^n)$．
  各 $x_i$ は可積分で，Fubini の定理と補題~\ref{lem:gauss-integral}より
  $\int x_i\,\d\gamma_n=\int t\varphi\,\d\lambda_1=0$，
  すなわち平均は $0$ である．
  共分散の $(i,j)$ 成分は $\int x_ix_j\,\d\gamma_n$ で，
  $i=j$ なら $\int t^2\varphi\,\d\lambda_1=1$，
  $i\neq j$ なら Fubini の定理で
  $\bigl(\int t\varphi\,\d\lambda_1\bigr)^2\cdot1=0$ となり，
  共分散は $I_n$ である．
\end{proof}

\begin{lemma}{直交不変性}{gauss-rotation}
  $U\in\R^{n\times n}$ が直交行列ならば
  $U\pushforward\gamma_n=\gamma_n$ である．
\end{lemma}

\begin{proof}
  $U^\top U=I$ より $(\det U)^2=1$，すなわち $\abs{\det U}=1$ である．
  任意の $B\in\Borel(\R^n)$ に対し，
  $\mathbf{1}_{U^{-1}B}(x)=\mathbf{1}_B(Ux)$ と
  アフィン変数変換公式（定理~\ref{thm:meas-affine-cov}，$f:=\mathbf{1}_B\varphi_n$）から
  \[
    \gamma_n(B)
    =\int\mathbf{1}_B(y)\,\varphi_n(y)\,\d\lambda_n(y)
    =\abs{\det U}\int\mathbf{1}_B(Ux)\,\varphi_n(Ux)\,\d\lambda_n(x)
    =\int\mathbf{1}_{U^{-1}B}(x)\,\varphi_n(x)\,\d\lambda_n(x),
  \]
  ここで $\norm{Ux}=\norm{x}$ より $\varphi_n(Ux)=\varphi_n(x)$ を用いた．
  右辺は $\gamma_n(U^{-1}B)=(U\pushforward\gamma_n)(B)$ である．
\end{proof}

\begin{lemma}{線形押し出しの標準形（鍵補題）}{gauss-linear-pf}
  任意の $A\in\R^{k\times n}$ に対し
  \[
    A\pushforward\gamma_n=(AA^\top)^{1/2}\pushforward\gamma_k.
  \]
\end{lemma}

\begin{proof}
  方針：SVD で $A$ を「直交 $\circ$ 座標のスケール $\circ$ 直交」に分解し，
  直交部分を直交不変性で吸収する．

  $r:=\rank A$ とし，SVD（補題~\ref{lem:la-svd}）で
  $A=\sum_{i=1}^rd_iu_iv_i^\top$ と書く．
  $r=0$（$A=O$）のときは両辺とも $\delta_0$（$0$ への押し出し）で等しいから，
  $r\geq1$ とする．
  正規直交系 $\{v_i\}_{i\leq r}$，$\{u_i\}_{i\leq r}$ をそれぞれ
  $\R^n$，$\R^k$ の正規直交基底 $\{v_i\}_{i=1}^n$，$\{u_i\}_{i=1}^k$ に
  延長し（主張~\ref{clm:la-onb-extension}），
  $V:=[v_1,\dots,v_n]$，$U:=[u_1,\dots,u_k]$ とおくと，
  列が正規直交基底なので $V,U$ は直交行列である．

  ここで写像 $D\colon\R^n\to\R^k$ と $D'\colon\R^k\to\R^k$ を
  \[
    D(z_1,\dots,z_n):=(d_1z_1,\dots,d_rz_r,0,\dots,0),
    \qquad
    D'(w_1,\dots,w_k):=(d_1w_1,\dots,d_rw_r,0,\dots,0)
  \]
  と定める（いずれも線形）．$V^\top x=(\inner{v_1}{x},\dots,\inner{v_n}{x})$ だから
  \[
    U\bigl(D(V^\top x)\bigr)
    =\sum_{i=1}^rd_i\inner{v_i}{x}\,u_i
    =Ax,
  \]
  すなわち $A=U\circ D\circ V^\top$ である．
  また $B:=\sum_{i=1}^rd_iu_iu_i^\top$ は対称で固有値 $d_i>0$ と $0$ をもつので
  半正定値，かつ $B^2=\sum_id_i^2u_iu_i^\top=AA^\top$
  （補題~\ref{lem:la-svd}）だから，平方根の一意性
  （定理~\ref{thm:la-sqrt}）より $(AA^\top)^{1/2}=B$ であり，
  同じ計算で $B=U\circ D'\circ U^\top$ である．

  押し出しの合成則（主張~\ref{clm:meas-pf-compose}）と
  直交不変性（補題~\ref{lem:gauss-rotation}；$V^\top,U^\top$ も直交行列）より
  \[
    A\pushforward\gamma_n
    =U\pushforward\bigl(D\pushforward({V^\top}\pushforward\gamma_n)\bigr)
    =U\pushforward(D\pushforward\gamma_n),
    \qquad
    B\pushforward\gamma_k
    =U\pushforward(D'\pushforward\gamma_k)
  \]
  だから，$D\pushforward\gamma_n=D'\pushforward\gamma_k$ を示せばよい．
  $\iota\colon\R^r\to\R^k$，$\iota(z):=(z_1,\dots,z_r,0,\dots,0)$ と
  $D_r:=\diag(d_1,\dots,d_r)\colon\R^r\to\R^r$ を使うと
  $D=\iota\circ D_r\circ\mathrm{pr}_r$，
  $D'=\iota\circ D_r\circ\mathrm{pr}_r'$
  （$\mathrm{pr}_r,\mathrm{pr}_r'$ はそれぞれ $\R^n,\R^k$ から
  最初の $r$ 座標への射影）と書ける．
  主張~\ref{clm:gauss-standard-moments}(i) より
  ${\mathrm{pr}_r}\pushforward\gamma_n=\gamma_r={\mathrm{pr}_r'}\pushforward\gamma_k$
  だから，合成則より
  \[
    D\pushforward\gamma_n
    =\iota\pushforward\bigl({D_r}\pushforward\gamma_r\bigr)
    =D'\pushforward\gamma_k.  \qedhere
  \]
\end{proof}

\begin{definition}{Gaussian 測度}{gauss-def}
  $m\in\R^n$ と $\Sigma\succeq0$ に対し，
  アフィン写像 $\ell(x):=m+\Sigma^{1/2}x$（連続ゆえ可測）による押し出し
  \[
    \Normal(m,\Sigma):=\ell\pushforward\gamma_n
  \]
  を，平均 $m$・共分散 $\Sigma$ の \textbf{Gaussian 測度}
  （正規分布）とよぶ
  （名称の正当化は主張~\ref{clm:gauss-moments}）．
  $\Sigma$ の正則性は仮定しない．
  例えば $\Sigma=O$ のとき $\ell\equiv m$ だから，
  $\Normal(m,O)$ は1点 $m$ に集中する\textbf{点測度}
  $\delta_m$（$\delta_m(B):=\mathbf{1}_B(m)$）である．
\end{definition}

\begin{lemma}{Gaussian 測度のアフィン像}{gauss-affine}
  $L\in\R^{k\times n}$，$b\in\R^k$ に対し
  \[
    (x\mapsto Lx+b)\pushforward\Normal(m,\Sigma)
    =\Normal(Lm+b,\;L\Sigma L^\top).
  \]
  （右辺は $L\Sigma L^\top\succeq0$（主張~\ref{clm:la-congruence}）
  より定義されている．）
\end{lemma}

\begin{proof}
  押し出しの合成則（主張~\ref{clm:meas-pf-compose}）より，左辺は
  合成写像 $x\mapsto L(m+\Sigma^{1/2}x)+b=(Lm+b)+(L\Sigma^{1/2})x$ による
  $\gamma_n$ の押し出しである．
  $A:=L\Sigma^{1/2}$ とおくと $AA^\top=L\Sigma L^\top$ だから，
  鍵補題（補題~\ref{lem:gauss-linear-pf}）と合成則より
  \begin{align*}
    (x\mapsto(Lm+b)+Ax)\pushforward\gamma_n
    &=(y\mapsto(Lm+b)+y)\pushforward\bigl(A\pushforward\gamma_n\bigr)\\
    &=(y\mapsto(Lm+b)+y)\pushforward
      \bigl((L\Sigma L^\top)^{1/2}\pushforward\gamma_k\bigr)\\
    &=\bigl(y\mapsto(Lm+b)+(L\Sigma L^\top)^{1/2}y\bigr)\pushforward\gamma_k
    =\Normal(Lm+b,\;L\Sigma L^\top).  \qedhere
  \end{align*}
\end{proof}

\begin{corollary}{結合 Gaussian 測度の周辺分布}{gauss-marginal}
  $m_1,m_2\in\R^n$，$\Sigma_1,\Sigma_2\succeq0$，$K\in\R^{n\times n}$ とし，
  $2n\times2n$ のブロック行列
  $A=\bigl(\begin{smallmatrix}\Sigma_1&K\\K^\top&\Sigma_2\end{smallmatrix}\bigr)$
  が $A\succeq0$ をみたすとする．
  このとき $\R^n\times\R^n=\R^{2n}$ 上の Gaussian 測度
  $\Normal\bigl((m_1,m_2),A\bigr)$ の第1周辺分布は $\Normal(m_1,\Sigma_1)$，
  第2周辺分布は $\Normal(m_2,\Sigma_2)$ である．
\end{corollary}

\begin{proof}
  第1射影 $\mathrm{pr}_1(x,y)=x$ は行列 $P_1:=[\,I\;\;O\,]\in\R^{n\times2n}$
  による線形写像であり，
  $P_1(m_1,m_2)=m_1$，$P_1AP_1^\top=\Sigma_1$（ブロック計算）だから，
  補題~\ref{lem:gauss-affine}より
  ${\mathrm{pr}_1}\pushforward\Normal((m_1,m_2),A)=\Normal(m_1,\Sigma_1)$．
  第2周辺も同様である．
\end{proof}

\begin{claim}{Gaussian 測度のモーメント}{gauss-moments}
  $\Normal(m,\Sigma)\in\Pp{2}(\R^n)$ であり，
  その平均ベクトルは $m$，共分散行列は $\Sigma$ である．
  特に $\int\norm{x}^2\,\d\Normal(m,\Sigma)=\norm{m}^2+\tr\Sigma$ であり，
  $\Wass_2$ の引数として意味をもつ（$\Pp{2}$ に属する）．
\end{claim}

\begin{proof}
  $\gamma_n\in\Pp{2}(\R^n)$ は平均 $0$・共分散 $I_n$ をもつ
  （主張~\ref{clm:gauss-standard-moments}）．
  $\Normal(m,\Sigma)=\ell\pushforward\gamma_n$
  （$\ell(x)=m+\Sigma^{1/2}x$）だから，
  アフィン像の平均・共分散（主張~\ref{clm:gauss-cov-basic}(iii)）より
  $\Normal(m,\Sigma)\in\Pp{2}(\R^n)$ で，平均は $m$，共分散は
  $\Sigma^{1/2}I_n\Sigma^{1/2}=\Sigma$ である．
  最後の等式は主張~\ref{clm:gauss-cov-basic}(ii) による．
\end{proof}

\begin{memo*}
  $\Sigma\succ0$（正則）のとき，$\Normal(m,\Sigma)$ は密度
  \[
    x\mapsto\det(2\pi\Sigma)^{-1/2}
    \exp\Bigl(-\tfrac{1}{2}(x-m)^\top\Sigma^{-1}(x-m)\Bigr)
  \]
  をもつ（アフィン変数変換公式から従う）．
  本資料ではこの表示は使わないため，証明は省略する．
  押し出しによる定義（定義~\ref{def:gauss-def}）の利点は，
  $\Sigma$ が特異で密度が存在しない場合も同じ定義がそのまま通ることにある．
\end{memo*}

\section{平均の分離}
\label{sec:gauss-centering}

\begin{lemma}{平均の分離}{gauss-centering}
  $\mu,\nu\in\Pp{2}(\R^n)$ とし，平行移動で中心化した測度を
  $\tilde\mu:=(x\mapsto x-m_\mu)\pushforward\mu$，
  $\tilde\nu:=(y\mapsto y-m_\nu)\pushforward\nu$ とおく
  （主張~\ref{clm:gauss-cov-basic}(iii) より
  $\tilde\mu,\tilde\nu\in\Pp{2}(\R^n)$ は平均 $0$ で，
  共分散は $\Sigma_\mu,\Sigma_\nu$ のまま変わらない）．このとき
  \[
    \Wass_2(\mu,\nu)^2
    =\norm{m_\mu-m_\nu}^2+\Wass_2(\tilde\mu,\tilde\nu)^2.
  \]
  この等式は Gaussian に限らず $\Pp{2}(\R^n)$ の任意の測度で成り立つ．
\end{lemma}

\begin{proof}
  $T(x,y):=(x-m_\mu,\,y-m_\nu)$ とおく（アフィンゆえ可測）．
  $\pi\in\Couplings(\mu,\nu)$ に対し
  $T\pushforward\pi\in\Couplings(\tilde\mu,\tilde\nu)$ である：実際
  \[
    (T\pushforward\pi)(A\times\R^n)
    =\pi\bigl((A+m_\mu)\times\R^n\bigr)
    =\mu(A+m_\mu)
    =\tilde\mu(A)
  \]
  （$A+m_\mu:=\{a+m_\mu\mid a\in A\}$），第2周辺も同様だからである．
  逆向きの平行移動 $T^{-1}(x,y)=(x+m_\mu,\,y+m_\nu)$ も同じ形なので，
  合成則（主張~\ref{clm:meas-pf-compose}）より対応
  $\pi\mapsto T\pushforward\pi$ は
  $\Couplings(\mu,\nu)$ と $\Couplings(\tilde\mu,\tilde\nu)$ の間の
  全単射である．

  コストを比べる．$\pi\in\Couplings(\mu,\nu)$ 上で
  $\tilde x:=x-m_\mu$，$\tilde y:=y-m_\nu$ と略記すると
  \[
    \norm{x-y}^2
    =\norm{\tilde x-\tilde y}^2
    +2\inner{\tilde x-\tilde y}{\,m_\mu-m_\nu}
    +\norm{m_\mu-m_\nu}^2.
  \]
  ここで周辺積分公式（補題~\ref{lem:meas-marginal-integral}）より
  $\int\tilde x\,\d\pi=\int(x-m_\mu)\,\d\mu=0$，
  $\int\tilde y\,\d\pi=0$（成分ごとに可積分：定義~\ref{def:gauss-mean-cov}）
  だから，交差項は積分すると消える．また変数変換公式
  （補題~\ref{lem:meas-change-of-variables}）より
  $\int\norm{\tilde x-\tilde y}^2\,\d\pi
  =\int\norm{x'-y'}^2\,\d(T\pushforward\pi)$ である．よって
  \[
    \int\norm{x-y}^2\,\d\pi
    =\int\norm{x'-y'}^2\,\d(T\pushforward\pi)
    +\norm{m_\mu-m_\nu}^2.
  \]
  対応 $\pi\mapsto T\pushforward\pi$ は全単射だから，
  両辺の下限を取れば主張の等式を得る．
\end{proof}

\section{Coupling のコストと相互共分散}
\label{sec:gauss-cost}

\begin{definition}{相互共分散行列}{gauss-cross-cov}
  $\mu,\nu\in\Pp{2}(\R^n)$ をともに平均 $0$ とし，
  $\pi\in\Couplings(\mu,\nu)$ とする．$\pi$ の\textbf{相互共分散行列}を
  \[
    K_\pi:=\int_{\R^n\times\R^n}xy^\top\,\d\pi(x,y)\in\R^{n\times n}
  \]
  と成分ごとの積分で定める．
  可積分性は $\abs{x_iy_j}\leq\frac{1}{2}(\norm{x}^2+\norm{y}^2)$ と
  周辺積分公式（補題~\ref{lem:meas-marginal-integral}）から従う．
\end{definition}

\begin{lemma}{coupling のコスト公式}{gauss-coupling-cost}
  上の設定で，$\mu,\nu$ の共分散行列を $\Sigma_1,\Sigma_2$ とすると，
  任意の $\pi\in\Couplings(\mu,\nu)$ に対して次が成り立つ．
  \begin{enumerate}
    \item $\displaystyle\int\norm{x-y}^2\,\d\pi
      =\tr\Sigma_1+\tr\Sigma_2-2\tr K_\pi$．
    \item ブロック行列
      $\begin{pmatrix}\Sigma_1&K_\pi\\K_\pi^\top&\Sigma_2\end{pmatrix}$
      は半正定値である．
  \end{enumerate}
\end{lemma}

\begin{proof}
  (i)：$\norm{x-y}^2=\norm{x}^2+\norm{y}^2-2\inner{x}{y}$ を積分する．
  周辺積分公式（補題~\ref{lem:meas-marginal-integral}）と
  主張~\ref{clm:gauss-cov-basic}(ii)（平均 $0$）より
  $\int\norm{x}^2\,\d\pi=\int\norm{x}^2\,\d\mu=\tr\Sigma_1$，
  同様に $\int\norm{y}^2\,\d\pi=\tr\Sigma_2$，
  さらに $\int\inner{x}{y}\,\d\pi=\sum_i\int x_iy_i\,\d\pi=\tr K_\pi$ である．

  (ii)：任意の $u,v\in\R^n$ に対し，
  積分の線形性（成分ごと）と周辺積分公式から
  \[
    \begin{pmatrix}u\\v\end{pmatrix}^{\!\top}
    \begin{pmatrix}\Sigma_1&K_\pi\\K_\pi^\top&\Sigma_2\end{pmatrix}
    \begin{pmatrix}u\\v\end{pmatrix}
    =u^\top\Sigma_1u+2\,u^\top K_\pi v+v^\top\Sigma_2v
    =\int\bigl(\inner{u}{x}+\inner{v}{y}\bigr)^2\,\d\pi
    \geq0.  \qedhere
  \]
\end{proof}

\section{トレース最大化と下界}
\label{sec:gauss-lower-bound}

次の命題が本章の心臓部である．
補題~\ref{lem:gauss-coupling-cost}により
「coupling のコストを最小化する」問題は
「ブロック半正定値制約の下で $\tr K$ を最大化する」
純粋な線形代数の問題に帰着する．

\begin{proposition}{トレース最大化}{gauss-trace-max}
  $\Sigma_1,\Sigma_2\succeq0$ に対し
  \[
    \max\left\{\tr K
    \;\middle|\;
    K\in\R^{n\times n},\;
    \begin{pmatrix}\Sigma_1&K\\K^\top&\Sigma_2\end{pmatrix}\succeq0
    \right\}
    =\tr\bigl(\Sigma_1^{1/2}\Sigma_2\Sigma_1^{1/2}\bigr)^{1/2}
  \]
  であり，最大値は実際に達成される．
\end{proposition}

\begin{proof}
  $N:=\Sigma_2^{1/2}\Sigma_1^{1/2}$ とおくと
  \[
    N^\top N
    =\Sigma_1^{1/2}\Sigma_2^{1/2}\Sigma_2^{1/2}\Sigma_1^{1/2}
    =\Sigma_1^{1/2}\Sigma_2\Sigma_1^{1/2}
  \]
  だから，特異値の和とトレースの関係（補題~\ref{lem:la-trace-norm}）より
  \[
    \tr\bigl(\Sigma_1^{1/2}\Sigma_2\Sigma_1^{1/2}\bigr)^{1/2}
    =\sum_i\sigma_i(N).
  \]

  \textbf{上界．}
  $K$ を制約をみたす任意の行列とすると，
  ブロック半正定値行列の特徴づけ（補題~\ref{lem:la-block-psd}）より
  $K=\Sigma_1^{1/2}C\Sigma_2^{1/2}$（$\norm{C}_{\mathrm{op}}\leq1$）と書ける．
  トレースの巡回性（主張~\ref{clm:la-trace-basic}）と
  von Neumann 型トレース不等式（補題~\ref{lem:la-vn-trace}）から
  \[
    \tr K
    =\tr\bigl(C\,\Sigma_2^{1/2}\Sigma_1^{1/2}\bigr)
    =\tr(CN)
    \leq\sum_i\sigma_i(N).
  \]

  \textbf{達成．}
  SVD（補題~\ref{lem:la-svd}）で $N=\sum_{i=1}^rd_iu_iv_i^\top$ と書き，
  $C^*:=\sum_{i=1}^rv_iu_i^\top$，
  $K^*:=\Sigma_1^{1/2}C^*\Sigma_2^{1/2}$ とおく．
  Bessel の不等式（主張~\ref{clm:la-bessel}）より，任意の $w$ で
  \[
    \norm{C^*w}^2
    =\Bigl\lVert\sum_i\inner{u_i}{w}\,v_i\Bigr\rVert^2
    =\sum_i\inner{u_i}{w}^2
    \leq\norm{w}^2
  \]
  だから $\norm{C^*}_{\mathrm{op}}\leq1$ であり，
  補題~\ref{lem:la-block-psd}より $K^*$ は制約をみたす．そのコストは
  \[
    \tr K^*
    =\tr(C^*N)
    =\tr\Bigl(\sum_{i,j}d_j\,v_iu_i^\top u_jv_j^\top\Bigr)
    =\sum_{i}d_i\tr(v_iv_i^\top)
    =\sum_id_i
    =\sum_i\sigma_i(N)
  \]
  であり，上界に一致する．
\end{proof}

\begin{proposition}{下界：同じ平均・共分散をもつ任意の測度}{gauss-lower-bound}
  $\mu,\nu\in\Pp{2}(\R^n)$ の平均を $m_1,m_2$，共分散を
  $\Sigma_1,\Sigma_2$ とすると
  \[
    \Wass_2(\mu,\nu)^2
    \geq
    \norm{m_1-m_2}^2
    +\tr\Sigma_1+\tr\Sigma_2
    -2\tr\bigl(\Sigma_1^{1/2}\Sigma_2\Sigma_1^{1/2}\bigr)^{1/2}.
  \]
  証明のどこにも Gaussian 性は使われないことに注意する．
\end{proposition}

\begin{proof}
  平均の分離（補題~\ref{lem:gauss-centering}）より，
  平均 $0$・共分散 $\Sigma_1,\Sigma_2$ の場合に
  $\Wass_2(\tilde\mu,\tilde\nu)^2\geq
  \tr\Sigma_1+\tr\Sigma_2-2\tr(\Sigma_1^{1/2}\Sigma_2\Sigma_1^{1/2})^{1/2}$
  を示せばよい．
  任意の $\pi\in\Couplings(\tilde\mu,\tilde\nu)$ に対し，
  コスト公式（補題~\ref{lem:gauss-coupling-cost}）より
  $K_\pi$ はトレース最大化（命題~\ref{prop:gauss-trace-max}）の
  制約をみたし，
  \[
    \int\norm{x-y}^2\,\d\pi
    =\tr\Sigma_1+\tr\Sigma_2-2\tr K_\pi
    \geq
    \tr\Sigma_1+\tr\Sigma_2
    -2\tr\bigl(\Sigma_1^{1/2}\Sigma_2\Sigma_1^{1/2}\bigr)^{1/2}.
  \]
  $\pi$ について下限を取れば結論を得る．
\end{proof}

\section{主定理}
\label{sec:gauss-main}

\begin{proposition}{結合 Gaussian による達成}{gauss-attainment}
  $m_1,m_2\in\R^n$，$\Sigma_1,\Sigma_2\succeq0$ とし，
  $K^*$ をトレース最大化（命題~\ref{prop:gauss-trace-max}）の最大化元，
  $A^*:=\bigl(\begin{smallmatrix}\Sigma_1&K^*\\K^{*\top}&\Sigma_2\end{smallmatrix}\bigr)\succeq0$
  とおく．このとき
  \[
    \gamma^*:=\Normal\bigl((m_1,m_2),\,A^*\bigr)
    \in\Couplings\bigl(\Normal(m_1,\Sigma_1),\,\Normal(m_2,\Sigma_2)\bigr)
  \]
  であり，そのコストは
  \[
    \int\norm{x-y}^2\,\d\gamma^*
    =\norm{m_1-m_2}^2
    +\tr\Sigma_1+\tr\Sigma_2-2\tr K^*.
  \]
\end{proposition}

\begin{proof}
  結合 Gaussian の周辺分布（系~\ref{cor:gauss-marginal}）より
  $\gamma^*$ の周辺分布は $\Normal(m_1,\Sigma_1)$，$\Normal(m_2,\Sigma_2)$
  だから $\gamma^*$ は coupling である．
  コストを計算する．$L:=[\,I\;\;-I\,]\in\R^{n\times2n}$ とおくと，
  $w=(x,y)$ に対し $x-y=Lw$ だから，
  変数変換公式（補題~\ref{lem:meas-change-of-variables}）と
  アフィン像（補題~\ref{lem:gauss-affine}）より
  \[
    \int\norm{x-y}^2\,\d\gamma^*
    =\int\norm{z}^2\,\d\bigl(L\pushforward\gamma^*\bigr)(z),
    \qquad
    L\pushforward\gamma^*
    =\Normal\bigl(m_1-m_2,\;LA^*L^\top\bigr).
  \]
  Gaussian 測度のモーメント（主張~\ref{clm:gauss-moments}）より
  \[
    \int\norm{z}^2\,\d\Normal\bigl(m_1-m_2,\,LA^*L^\top\bigr)
    =\norm{m_1-m_2}^2+\tr\bigl(LA^*L^\top\bigr)
  \]
  であり，ブロック計算で
  $LA^*L^\top=\Sigma_1-K^*-K^{*\top}+\Sigma_2$，
  $\tr K^{*\top}=\tr K^*$ だから
  $\tr(LA^*L^\top)=\tr\Sigma_1+\tr\Sigma_2-2\tr K^*$ である．
\end{proof}

\begin{theorem}{Gaussian 測度の間の $W_2$ 距離（Givens--Shortt）}{gauss-main}
  $m_1,m_2\in\R^n$，$\Sigma_1,\Sigma_2\succeq0$（特異でもよい）に対し
  \[
    \Wass_2\bigl(\Normal(m_1,\Sigma_1),\,\Normal(m_2,\Sigma_2)\bigr)^2
    =
    \norm{m_1-m_2}^2
    +\tr\Bigl(\Sigma_1+\Sigma_2
      -2\bigl(\Sigma_1^{1/2}\Sigma_2\Sigma_1^{1/2}\bigr)^{1/2}\Bigr)
  \]
  が成り立ち，下限は結合 Gaussian $\gamma^*$
  （命題~\ref{prop:gauss-attainment}）で達成される（$\inf$ は $\min$ である）．
\end{theorem}

\begin{proof}
  主張~\ref{clm:gauss-moments}より $\Normal(m_i,\Sigma_i)\in\Pp{2}(\R^n)$ で，
  その平均は $m_i$，共分散は $\Sigma_i$ である．
  よって下界命題~\ref{prop:gauss-lower-bound}から
  $\Wass_2^2\geq(\text{右辺})$．
  逆に，命題~\ref{prop:gauss-attainment}の coupling $\gamma^*$ のコストは，
  $\tr K^*=\tr(\Sigma_1^{1/2}\Sigma_2\Sigma_1^{1/2})^{1/2}$
  （命題~\ref{prop:gauss-trace-max}）より右辺に等しい．
  したがって $\Wass_2^2\leq(\text{右辺})$ であり，等号と達成が従う．
\end{proof}

\begin{remark}{公式の対称性}{gauss-symmetry}
  $\Wass_2$ は距離だから右辺も $\Sigma_1\leftrightarrow\Sigma_2$ の
  交換で不変なはずである．実際，
  $N=\Sigma_2^{1/2}\Sigma_1^{1/2}$ に対し
  $\sigma_i(N)=\sigma_i(N^\top)$（定義~\ref{def:la-singular}）と
  補題~\ref{lem:la-trace-norm}から
  \[
    \tr\bigl(\Sigma_1^{1/2}\Sigma_2\Sigma_1^{1/2}\bigr)^{1/2}
    =\sum_i\sigma_i(N)
    =\sum_i\sigma_i(N^\top)
    =\tr\bigl(\Sigma_2^{1/2}\Sigma_1\Sigma_2^{1/2}\bigr)^{1/2}
  \]
  であり，確かに対称である（原論文 p.~239 の
  $\tr\sqrt{B^\top B}=\tr\sqrt{BB^\top}$ による注意に対応）．
\end{remark}

\begin{remark}{最適輸送写像}{gauss-optimal-map}
  $\Sigma_1\succ0$ のときは，最適な輸送を\textbf{写像}で書ける．
  $M:=\Sigma_1^{1/2}\Sigma_2\Sigma_1^{1/2}$，
  $T:=\Sigma_1^{-1/2}M^{1/2}\Sigma_1^{-1/2}$
  （合同変換なので $T\succeq0$：主張~\ref{clm:la-congruence}），
  $\phi(x):=m_2+T(x-m_1)$ とおくと：
  \begin{itemize}
    \item $T\Sigma_1T
      =\Sigma_1^{-1/2}M^{1/2}\bigl(\Sigma_1^{-1/2}\Sigma_1\Sigma_1^{-1/2}\bigr)
        M^{1/2}\Sigma_1^{-1/2}
      =\Sigma_1^{-1/2}M\Sigma_1^{-1/2}=\Sigma_2$ だから，
      アフィン像（補題~\ref{lem:gauss-affine}）より
      $\phi\pushforward\Normal(m_1,\Sigma_1)=\Normal(m_2,\Sigma_2)$．
    \item $(\Id,\phi)(x):=(x,\phi(x))$ もアフィン写像で，
      $\pi_\phi:=(\Id,\phi)\pushforward\Normal(m_1,\Sigma_1)$ は
      coupling になる（$\mathrm{pr}_1\circ(\Id,\phi)=\Id$，
      $\mathrm{pr}_2\circ(\Id,\phi)=\phi$ と合成則）．
    \item そのコストは，$x-\phi(x)=(I-T)(x-m_1)+(m_1-m_2)$ と
      主張~\ref{clm:gauss-cov-basic}(ii)，交差項の消滅
      （$\int(x-m_1)\,\d\Normal(m_1,\Sigma_1)=0$）より
      \[
        \int\norm{x-\phi(x)}^2\,\d\Normal(m_1,\Sigma_1)
        =\norm{m_1-m_2}^2+\tr\bigl((I-T)\Sigma_1(I-T)\bigr).
      \]
      展開すると
      $(I-T)\Sigma_1(I-T)=\Sigma_1+T\Sigma_1T-T\Sigma_1-\Sigma_1T$ で，
      巡回性から
      $\tr(T\Sigma_1)=\tr(\Sigma_1^{-1/2}M^{1/2}\Sigma_1^{1/2})
      =\tr(M^{1/2}\Sigma_1^{1/2}\Sigma_1^{-1/2})=\tr M^{1/2}$
      （$\tr(\Sigma_1T)$ も同じ）だから，コストは
      $\norm{m_1-m_2}^2+\tr(\Sigma_1+\Sigma_2)-2\tr M^{1/2}
      =\Wass_2^2$ に一致する．
  \end{itemize}
  すなわち $\pi_\phi$ も最適 coupling であり，
  Gaussian の場合は最適輸送が決定論的な（アフィン写像による）
  輸送で実現できる．
\end{remark}

\begin{memo*}
  $\Sigma_1\succ0$ のとき，最適 coupling は上の $\pi_\phi$ に
  限ることが知られている（Knott--Smith，Brenier の定理）．
  証明には Kantorovich 双対性など本資料の範囲を超える道具が要るため，
  Villani (2009) Chapter 2，Peyré--Cuturi (2019) の 2.6 節に譲る．
\end{memo*}

\begin{remark}{原論文との対応}{gauss-gs-correspondence}
  定理~\ref{thm:gauss-main}は Givens--Shortt (1984) の
  Proposition 7（p.~236，式 (6)）である．ただし：
  \begin{itemize}
    \item 原論文は $\Sigma_1,\Sigma_2$ が正則な場合を証明し，
      特異な場合は「同様の議論が通るが煩雑な場合分けを要する」
      （p.~239）と注記するに留まる．
      本章の証明は Schur 補行列（$\Sigma_1^{-1}$ を要する）の代わりに
      補題~\ref{lem:la-block-psd}を使うため，特異な場合も区別なく扱えた．
    \item 原論文の Lemma 2（最適 coupling は Gaussian に取れる）は，
      最適 coupling の存在（同論文 Proposition 1：弱コンパクト性）に
      依存する．本章では下界（命題~\ref{prop:gauss-lower-bound}）と
      達成（命題~\ref{prop:gauss-attainment}）を分けて示したので，
      存在定理は不要になった．
      「任意の coupling の相互共分散はブロック半正定値制約をみたす」
      （下界方向）と
      「制約をみたす任意の行列は結合 Gaussian で実現できる」
      （達成方向）は別の主張であり，Gaussian 性を使うのは後者だけである．
    \item 原論文のトレース最大化は Schur 補行列によるパラメータ付け・
      Stiefel 制約下の Lagrange 乗数法・Courant--Fischer の定理
      （pp.~237--239）による．
      本章では von Neumann 型トレース不等式（補題~\ref{lem:la-vn-trace}）で
      置き換えた．
    \item 先行研究 Dowson--Landau
      （\emph{The Fréchet distance between multivariate normal
      distributions}, J.\ Multivariate Anal.\ \textbf{12} (1982), 450--455）の
      公式は共分散が可換な場合のみ正しい（原論文 p.~240 の謝辞）．
      可換な場合は系~\ref{cor:gauss-commuting}を参照．
  \end{itemize}
\end{remark}

\section{系と例}
\label{sec:gauss-corollaries}

\begin{corollary}{Gelbrich の下界}{gauss-gelbrich}
  任意の $\mu,\nu\in\Pp{2}(\R^n)$ に対し
  \[
    \Wass_2(\mu,\nu)
    \geq
    \Wass_2\bigl(\Normal(m_\mu,\Sigma_\mu),\,\Normal(m_\nu,\Sigma_\nu)\bigr).
  \]
  すなわち，同じ平均・共分散をもつ測度の組の中で，
  Gaussian の組は $\Wass_2$ を最小にする．
\end{corollary}

\begin{proof}
  下界命題~\ref{prop:gauss-lower-bound}の右辺は，
  主定理~\ref{thm:gauss-main}より
  $\Wass_2(\Normal(m_\mu,\Sigma_\mu),\Normal(m_\nu,\Sigma_\nu))^2$
  に等しい．
\end{proof}

\begin{remark}{Gelbrich の下界について}{gauss-gelbrich-note}
  系~\ref{cor:gauss-gelbrich}は Gelbrich
  （\emph{On a formula for the $L^2$ Wasserstein metric between measures on Euclidean and Hilbert spaces},
  Math.\ Nachr.\ \textbf{147} (1990), 185--203）による．
  Givens--Shortt の原論文には現れないが，
  本章の証明構成（下界が Gaussian 性を使わないこと）からは
  追加の労力なしに得られる．
\end{remark}

\begin{corollary}{1次元の場合}{gauss-1d}
  $m_1,m_2\in\R$，$\sigma_1,\sigma_2\geq0$ に対し
  \[
    \Wass_2\bigl(\Normal(m_1,\sigma_1^2),\,\Normal(m_2,\sigma_2^2)\bigr)^2
    =(m_1-m_2)^2+(\sigma_1-\sigma_2)^2.
  \]
\end{corollary}

\begin{proof}
  $n=1$ では $\Sigma_i=\sigma_i^2\geq0$ はスカラーで，
  $\bigl(\sigma_1\sigma_2^2\sigma_1\bigr)^{1/2}=\sigma_1\sigma_2$
  （$\sigma_i\geq0$ に注意）だから，主定理~\ref{thm:gauss-main}より
  \[
    \Wass_2^2
    =(m_1-m_2)^2+\sigma_1^2+\sigma_2^2-2\sigma_1\sigma_2
    =(m_1-m_2)^2+(\sigma_1-\sigma_2)^2.  \qedhere
  \]
\end{proof}

\begin{corollary}{可換な共分散の場合}{gauss-commuting}
  $\Sigma_1\Sigma_2=\Sigma_2\Sigma_1$ ならば
  \[
    \Wass_2\bigl(\Normal(m_1,\Sigma_1),\,\Normal(m_2,\Sigma_2)\bigr)^2
    =\norm{m_1-m_2}^2
    +\tr\bigl(\bigl(\Sigma_1^{1/2}-\Sigma_2^{1/2}\bigr)^2\bigr).
  \]
\end{corollary}

\begin{proof}
  可換なら $(\Sigma_1^{1/2}\Sigma_2\Sigma_1^{1/2})^{1/2}
  =\Sigma_1^{1/2}\Sigma_2^{1/2}$（系~\ref{cor:la-commuting-sqrt}）．
  トレースの線形性と
  $\bigl(\Sigma_1^{1/2}-\Sigma_2^{1/2}\bigr)^2
  =\Sigma_1+\Sigma_2-\Sigma_1^{1/2}\Sigma_2^{1/2}-\Sigma_2^{1/2}\Sigma_1^{1/2}$，
  および巡回性
  $\tr(\Sigma_2^{1/2}\Sigma_1^{1/2})=\tr(\Sigma_1^{1/2}\Sigma_2^{1/2})$ から
  \[
    \tr\Bigl(\Sigma_1+\Sigma_2-2\Sigma_1^{1/2}\Sigma_2^{1/2}\Bigr)
    =\tr\bigl(\bigl(\Sigma_1^{1/2}-\Sigma_2^{1/2}\bigr)^2\bigr).  \qedhere
  \]
\end{proof}

\begin{example}{点測度：特異な場合の検算}{gauss-dirac}
  $\Sigma_1=\Sigma_2=O$ のとき $\Normal(m_i,O)=\delta_{m_i}$
  （定義~\ref{def:gauss-def}）で，主定理より
  $\Wass_2(\delta_{m_1},\delta_{m_2})=\norm{m_1-m_2}$ となる．
  これは直接の計算とも合う：
  $\Couplings(\delta_{m_1},\delta_{m_2})=\{\delta_{(m_1,m_2)}\}$
  （coupling の全測度は $\{m_1\}\times\{m_2\}$ に集中する）だから
  コストは常に $\norm{m_1-m_2}^2$ である．
  共分散が特異な場合を許した定式化の恩恵が最も端的に現れる例である．
\end{example}

\begin{example}{等方的な場合}{gauss-isotropic}
  $\Sigma_i=\sigma_i^2I_n$（$\sigma_i\geq0$）は可換だから，
  系~\ref{cor:gauss-commuting}より
  \[
    \Wass_2^2
    =\norm{m_1-m_2}^2+\tr\bigl((\sigma_1-\sigma_2)^2I_n\bigr)
    =\norm{m_1-m_2}^2+n\,(\sigma_1-\sigma_2)^2.
  \]
\end{example}

\begin{memo*}
  平均 $0$ の場合の右辺
  $\bigl(\tr(\Sigma_1+\Sigma_2-2(\Sigma_1^{1/2}\Sigma_2\Sigma_1^{1/2})^{1/2})\bigr)^{1/2}$
  は，量子情報理論で密度行列の間の距離として知られる
  \textbf{Bures 距離}と（正規化を除いて）一致し，
  \textbf{Bures--Wasserstein 距離}ともよばれる
  （Bhatia--Jain--Lim,
  \emph{On the Bures--Wasserstein distance between positive definite matrices},
  Expo.\ Math.\ \textbf{37} (2019), 165--191）．
  また原論文の COROLLARY（p.~239）は $n=2$ の行列式による表示
  \[
    \Wass_2\bigl(\Normal(0,\Sigma_1),\Normal(0,\Sigma_2)\bigr)
    =\sqrt{\tr\Sigma_1+\tr\Sigma_2
      -2\bigl[\tr(\Sigma_1\Sigma_2)+2\sqrt{\det(\Sigma_1\Sigma_2)}\bigr]^{1/2}}
  \]
  を与える（Cayley--Hamilton の定理による．証明は原論文参照）．
\end{memo*}
